% Generated draft: requires mathematical and source review.
\documentclass[11pt]{article}
\usepackage[margin=25mm]{geometry}
\usepackage{fontspec}
\setmainfont{Latin Modern Roman}
\usepackage{amsmath,amssymb,mathtools,unicode-math}
\setmathfont{Latin Modern Math}
\usepackage{enumitem,needspace,xurl,hyperref}
\setlength{\parindent}{0pt}
\setlength{\parskip}{5pt}
\setlength{\emergencystretch}{3em}
\newcommand{\N}{\mathbb{N}}
\newcommand{\Z}{\mathbb{Z}}
\newcommand{\Q}{\mathbb{Q}}
\newcommand{\R}{\mathbb{R}}
\newcommand{\C}{\mathbb{C}}
\newcommand{\F}{\mathbb{F}}
\newcommand{\A}{\mathbb{A}}
\newcommand{\PP}{\mathbb{P}}
\newcommand{\E}{\mathbb{E}}
\newcommand{\eps}{\varepsilon}
\newcommand{\dd}{\,\mathrm{d}}
\begin{document}
\section*{UnsolvedMath ID: 800942}
\textbf{Source problem code:} catalogue entry
\subsection{Definitions and mathematical statement}
A closed Riemannian manifold is a compact Riemannian manifold without boundary. Let $(M,g)$ be such a manifold with negative sectional curvature, meaning that $K_g(\sigma)<0$ for every tangent two-plane $\sigma$. Let $N(T)$ denote the source's number of closed geodesics of length at most $T$. The source does not specify whether geodesics must be primitive, whether iterates are counted separately, or whether opposite orientations are identified; these conventions remain unresolved.

Write $h>0$ for the constant appearing in the asymptotic quoted by the source,
\[
N(T)\sim \frac{e^{hT}}{hT}\qquad(T\to\infty).
\]
The supplied excerpt does not define $h$ independently. For $x>2$, define
\[
L(x):=\int_2^x\frac{\dd u}{\log u}.
\]
The intended assertion of Conjecture 22 is that, for every $(M,g)$ as above, there exists $\varepsilon>0$ such that
\[
N(T)=\bigl(1+O(e^{-\varepsilon T})\bigr)L(e^{hT})
\qquad(T\to\infty).
\]
Explicitly, there should exist constants $\varepsilon>0$, $C>0$, and $T_0>\max\{0,(\log 2)/h\}$, allowed to depend on $(M,g)$, such that for every real $T\ge T_0$,
\[
\left|N(T)-L(e^{hT})\right|
\le C e^{-\varepsilon T}L(e^{hT}).
\]
This remains a qualified formulation: the source does not specify connectedness or a dimension restriction, and the counting convention and meaning of $h$ require confirmation. In particular, the displayed leading asymptotic cannot be extended without qualification to arbitrary disconnected unions: two isometric components double the counting function.
\subsection{Short English statement}
The conjecture asks whether the closed-geodesic counting function on every closed negatively curved Riemannian manifold has an exponentially small relative error compared with $\int_2^{e^{hT}}\dd u/\log u$. The counting conventions, the independent definition of $h$, and the intended connectedness and dimension assumptions require clarification.
\subsection{Sources}
\begin{itemize}
\item Local UnsolvedMath registry: permanent ID $800942$, source problem code \texttt{catalogue entry}, AMR Open Problem Lists. The supplied registry attributes the problem to Burns and Matveev, gives proposed year 2019, and labels it \emph{open}; that status is supplied metadata, not an independently verified assessment.
\item Burns and Matveev, \emph{Open Problems and Questions about Geodesics} (2019, as identified in the supplied source list), Conjecture 22. Supplied source URL: \url{https://arxiv.org/abs/1308.5417}. The background records access on 2026-10-16 and extraction by \texttt{arxiv\_latex}; these are supplied provenance notes, not a claim of independent access or verification.
\item The supplied statement attributes the leading asymptotic $N(T)\sim e^{hT}/(hT)$ to the reference label \texttt{[Margulis2004]}. No full bibliographic entry for that reference is supplied.
\end{itemize}
\end{document}
